\chapter{Происхождение весов из расширенной концевой бимодульной алгебры}
\label{ch:v8-baryon-c0-extended-endpoint-bimodule-weight-origin-gate}

Ограничимся пятью одномерными углами. Старую тройку составляют $X_L$,
$e_R$ и $X_R$; новую пару --- $s_0$ и $a_0$. До добавления старо-новых
стрелок локальная концевая алгебра имеет вид
\begin{equation}
 \mathcal A_{
m loc}
 =\mathbb C P_{X_L}\oplus\mathbb C P_{e_R}
 \oplus\mathbb C P_{X_R}\oplus M_2(\mathbb C)_{s_0a_0}.
 \label{eq:v8-baryon-c0-local-extended-endpoint-algebra}
\end{equation}
Её центральные действия на multiplicity-пространстве дают три точных
координатных проектора
\begin{equation}
 \Pi_0=\operatorname{diag}(1,0,0),\qquad
 \Pi_1=\operatorname{diag}(0,1,0),\qquad
 \Pi_2=\operatorname{diag}(0,0,1).
 \label{eq:v8-baryon-c0-bimodule-coordinate-projectors}
\end{equation}
Это устраняет недиагональное смешивание ветвей, но ещё не задаёт их порядок.

\section{Симплекс положительных следов}

Общий положительный нормированный след на локальной алгебре задаётся
\begin{equation}
 p_0,p_1,p_2,q>0,
 \qquad p_0+p_1+p_2+2q=1.
 \label{eq:v8-baryon-c0-local-trace-simplex}
\end{equation}
На трёх коннекторах он индуцирует диагональную форму
\begin{equation}
 K_{\rm conn}(p,q)
 =\operatorname{diag}(p_0+q,p_1+q,p_2+q).
 \label{eq:v8-baryon-c0-trace-induced-connector-weights}
\end{equation}
Два точных допустимых свидетеля
\begin{align}
 (p_0,p_1,p_2,q)_A&=(1/10,1/5,1/2,1/10),\nonumber\\
 (p_0,p_1,p_2,q)_B&=(1/2,1/10,1/5,1/10)
 \label{eq:v8-baryon-c0-two-local-trace-witnesses}
\end{align}
дают разные минимальные направления. Даже дополнительное вектороподобное
условие $p_0=p_2$ оставляет
\begin{equation}
 K_{\rm conn}=\operatorname{diag}(p_X+q,p_e+q,p_X+q),
 \label{eq:v8-baryon-c0-vectorlike-two-weight-residual}
\end{equation}
то есть свободное отношение $p_e/p_X$ и как минимум двукратную $X$-ветвь.

\section{Полное Hom-замыкание как условная альтернатива}

Если не выбирать одну стрелку, а включить внутренний мост $s_0-a_0$ и все
три допустимые старо-новые стрелки, граф углов содержит рёбра
\begin{equation}
 E_{\rm full}=\{s_0a_0,X_Ls_0,e_Ra_0,X_Ra_0\}.
 \label{eq:v8-baryon-c0-full-local-edge-set}
\end{equation}
Он связен, и произведения матричных единиц порождают
\begin{equation}
 \operatorname{Alg}(E_{\rm full})=M_5(\mathbb C),
 \qquad \dim_{\mathbb C}=25,
 \qquad Z=M_5'=\mathbb C I_5.
 \label{eq:v8-baryon-c0-full-local-m5-closure}
\end{equation}
Единственный нормированный след полного блока даёт
\begin{equation}
 \tau_5(P_j)=\frac15,
 \qquad K_{\rm conn}^{(5)}=\frac25 I_3.
 \label{eq:v8-baryon-c0-m5-unique-trace-connector-metric}
\end{equation}
После trace-ортонормировки всех трёх комплексных стрелок получаем шесть
вещественных квадратур:
\begin{equation}
 \mathcal F_{51}=\mathcal F_{45}\oplus
 \bigoplus_{j=0}^{2}\operatorname{span}_{\mathbb R}
 \{C_{j,1},C_{j,2}\},
 \qquad \rank\mathcal F_{51}=51.
 \label{eq:v8-baryon-c0-full-multiplicity-frame-51}
\end{equation}
Этот маршрут устраняет необходимость выбора точки $\mathbb{RP}^2$ для
марковской примитивности, но заменяет одиночный морфизм полным
ковариационным пакетом.

Итог различает два вопроса:
\begin{equation}
 \text{canonical full multiplicity frame: условно да},
 \qquad
 \text{single }c_0\text{ map: нет}.
 \label{eq:v8-baryon-c0-full-frame-versus-single-map-ledger}
\end{equation}

\begin{theorem}[Бимодульная диагонализация без весового выбора]
Расширенная концевая алгебра канонически различает три multiplicity-метки и
запрещает их недиагональное смешивание. Однако её семейство положительных
следов не выбирает относительные диагональные веса. Единственный след
возникает лишь после полного $M_5$-замыкания всеми тремя коннекторами.
\end{theorem}

\begin{nogocriterion}[Полный кадр не равен одиночной карте c0]
Нельзя использовать уникальность следа уже связного $M_5$ как вывод одной
из стрелок: сам полный блок был получен добавлением всех трёх стрелок.
Кадр $F_{51}$ является допустимой изотропной динамической архитектурой, но
не выбирает rank-one морфизм, требуемый исходной интерпретацией $c_0$.
\end{nogocriterion}

\begin{protocol}\begin{enumerate}
\item локальная алгебра до старо-новых стрелок: \textbf{$\mathbb C^3\oplus M_2$};
\item coordinate projectors: \textbf{$3$};
\item свободных локальных trace-весов после нормировки: \textbf{$3$};
\item анизотропный вес выбран: \textbf{нет};
\item полное Hom-замыкание: \textbf{$M_5$};
\item уникальный след полного блока: \textbf{да};
\item полный multiplicity-frame: \textbf{$6$ новых направлений, ранг $51$};
\item одиночная карта $c_0$ выбрана: \textbf{нет};
\item абсолютная скорость: \textbf{открыта};
\item следующий гейт: \textbf{совместимость полного multiplicity-frame с
одиночной картой $c_0$}.
\end{enumerate}\end{protocol}

Машинный сертификат сохранён в
\path{s2t_v8_baryon_c0_extended_endpoint_bimodule_weight_origin_gate_results.json}.